\setcounter{section}{-1}

\section{Prémisses et éléments de théorèmes}


\begin{definition}
Un \textbf{axiome} est un énoncé (le plus souvent évident) admis comme vrai par tous et ne nécessitant pas justification supplémentaire 
\end{definition}


\smallskip
\begin{exemples}{0mm}
\item Un segment de droite peut être tracé en joignant deux points distincts.
\item Par un point $P$ extérieure à une droite $d$ passe une et une seule droite parallèle à $d$.
\end{exemples}


\medskip
\begin{remarques}
\item Un énoncé est vrai s'il est toujours vrai.
\item Pour qu'un énoncé soit faux, il suffit qu'il le soit une seule fois. La recherche d'un contre-exemple est une bonne manière de démontrer qu'un énoncé est faux.
\end{remarques}


\smallskip
\begin{definition}
Une \textbs{conjecture} est un énoncé qui n'a pas (encore!) été démontré mais dont on a l'intuition qu'il est vrai.
\end{definition}

\smallskip
\begin{exemple}{0.9}\guillemets{Tout nombre entier pair strictement supérieur ou égal à 4 peut s'écrire comme la somme de deux nombres premiers.}\vspace{-3mm}

\small
\hfill \textit{Conjecture de Goldbach}
\end{exemple}

\bigskip
\begin{definition}
Un \textbf{théorème} est un énoncé qui a été démontré comme étant vrai. Il est constitué de 3 parties :

\begin{enumerate}[label=\arabic*),itemsep=2mm]
\item Hypothèse : \hspace{8mm}\begin{minipage}[t]{0.7\textwidth}
Pré-requis à l'application du théorème.

\smallskip
\small
Ce sont les conditions nécessaires qui doivent être vérifiées pour utiliser la conclusion du théorème. Les conditions sont des données connues et utilisables dans la démonstration du théorème
\end{minipage}
\item Conclusion :\hspace{7mm} \begin{minipage}[t]{0.7\textwidth}
Propriété applicable si les hypothèses sont vérifiées.
\end{minipage}
\item Démonstration :\hspace{1mm} \begin{minipage}[t]{0.7\textwidth}
Argumentation qui s'appuie sur des faits admis comme vrai par tous, qui ne peut pas être contredite et permettant d'établir la véracité d'une proposition mathématique.
\end{minipage}
\end{enumerate}

\bigskip
Un théorème dont la portée est moins importante est souvent appelé une \textbs{proposition}.
\end{definition}


\clearpage
La formulation mathématique habituelle d'un théorème  est de la forme :
\begin{center}
\textsl{\guillemets{\textbs{Si} hypothèse (H) \textbs{Alors} conclusion (C)}}

\medskip
ou encore

\medskip
\textsl{\guillemets{hypothèse (H) $\boldsymbol{\Rightarrow}$ conclusion (C)}}
\end{center}

Il faut remarquer que la négation d'un théorème, i.e.\footnote{du latin \guillemets{id est} signifiant \guillemets{c'est-à-dire}} prendre la négation de l'hypothèse et la négation de la conclusion, n'est pas nécessairement vraie. En d'autres termes :

$$ (H \Rightarrow C) \notimply (NH \Rightarrow NC)$$


\bigskip
A titre d'exemple voici un théorème bien connu et sa négation :

\medskip
\textbf{Théorème :} \guillemets{Si une fonction $f$ est dérivable en $a$ alors $f$ est continue en $a$.}

\medskip
\textbf{Négation :} \guillemets{Si une fonction $f$ n'est pas dérivable en $a$ alors $f$ n'est pas continue en $a$}

\bigskip
Si le théorème ci-dessus est bien entendu vrai, sa négation, elle, ne l'est pas!



\bigskip
\begin{enumerate}[label=\roman*),itemsep=0mm]
\item Donner un autre exemple de théorème où la négation de ce dernier est fausse.
\item Donner un exemple de théorème où la négation de ce dernier est vraie.
\end{enumerate}








\begin{definitions}
\begin{itemize}
\item Un \textbs{corollaire} est un théorème qui découle immédiatement et logiquement d'un autre théorème. Il ne nécessite généralement que très peu d'arguments pour le démontrer.
\item La \textbs{réciproque} d'un théorème est un énoncé qui s'obtient en inversant le rôle de l'hypothèse et de la conclusion.
\item La \textbs{contraposée} d'un théorème s'obtient en prenant la négation de la réciproque (hypothèse et conclusion).
\end{itemize}
\end{definitions}



\bigskip
\begin{remarques}
\item La réciproque d'un théorème n'est pas toujours vraie.
$$ (H \Rightarrow C) \notimply (C \Rightarrow H)$$
\item La contraposée d'un théorème est toujours vraie.
$$ (H \Rightarrow C) \imply (NC \Rightarrow NH)$$
\end{remarques}







%\begin{minipage}{0.38\textwidth}
%\psset{yunit=0.7cm,xunit=0.7cm,gridcolor=cgrille}
%\begin{pspicture*}(0,0)(10,4)
%\psgrid[subgriddiv=1,griddots=7,gridlabels=0pt](0,0)(10,4)
%\psaxes[Dy=1, Dx=1, labelFontSize=\scriptstyle]{->}(0,0)(10,4)
%\uput{0}[0]{0}(1,2){$(H\Rightarrow C)$}
%\uput{0}[0]{45}(4,2){\LARGE{$\cancel{\Rightarrow}$}}
%\uput{0}[0]{0}(5.5,3){$(C \Rightarrow H)$}
%\uput{0}[0]{-45}(4,2){\LARGE{$\Rightarrow$}}
%\uput{0}[0]{0}(5.5,1){$NC \Rightarrow NH$}
%\end{pspicture*}
%\end{minipage}












